\chapter{Introduction}
\label{chap:introduction}

\section{Background}

SVG is one of the most used graphics formats on the web.
It is expressive, resolution-independent, and naturally integrated with HTML/CSS workflows.
For icons, diagrams, charts, and UI assets, SVG is often the most practical authoring format.

The limitation appears in heavy scenes.
When documents contain many paths, groups, and style combinations, browser-side interpretation can become expensive.
At the same time, modern web applications increasingly depend on explicit GPU pipelines for predictable rendering cost.
This creates a gap: SVG is declarative and semantic, while GPU pipelines consume explicit primitives and buffers.

\texttt{Svg2GPU} addresses this gap by converting SVG input into typed, GPU-oriented data.
The goal is not a full browser-engine reimplementation, but a robust conversion pipeline that preserves relevant visual semantics and can be validated reproducibly.

\section{Problem statement}

The central problem of this thesis is:

\begin{quote}
How can SVG input be transformed into robust, render-ready geometry and style data for browser GPU workflows, while preserving correctness and keeping the system practical to maintain?
\end{quote}

This problem includes four technical challenges:

\begin{itemize}
\item reliable parsing of heterogeneous SVG primitives and attributes;
\item conversion from path semantics to geometric primitives usable by render APIs;
\item deterministic handling of style, winding, and transform semantics;
\item safe behavior for malformed or partially unsupported input.
\end{itemize}

\section{Objectives}

The project objectives are:

\begin{itemize}
\item define a typed internal model for supported SVG elements and styles;
\item implement parser and style normalization components;
\item implement geometry processing foundations (fill and stroke conversion flow);
\item integrate with browser rendering classes in a modular way;
\item provide executable validation for core functionalities F1 and F2;
\item document architecture and API behavior in a usable form.
\end{itemize}

\section{Method and approach}

Development follows a staged approach:

\begin{enumerate}
\item parse SVG source into typed objects;
\item normalize style information and validate required fields;
\item generate renderable geometry from parsed data;
\item submit geometry through render-layer abstractions;
\item verify behavior through automated tests and scenario-based inspection.
\end{enumerate}

The approach is intentionally incremental.
It prioritizes stable contracts and reproducible behavior before full feature coverage.
This is suitable for bachelor thesis scope and produces a foundation that can be extended after delivery.

\section{Research context and practical relevance}

The thesis is situated at the intersection of frontend engineering and geometry processing.
Many web teams already author graphics in SVG, but once product requirements demand real-time interaction or larger scene complexity, native SVG rendering is not always the most predictable option.
In these contexts, teams often move directly to low-level APIs and lose high-level semantic structure.

The practical relevance of \texttt{Svg2GPU} is to avoid this hard switch.
Instead of replacing SVG authoring, it introduces a transformation layer that preserves useful semantics while enabling explicit rendering workflows.
This makes the approach attractive for interactive editors, visualization tooling, and performance-sensitive web interfaces where both expressiveness and control are required.

\section{Scope}

\subsection{In scope}

\begin{itemize}
\item core SVG primitive parsing (path and common geometric elements);
\item style parsing and normalization for practical rendering needs;
\item path-focused geometry conversion workflow;
\item practical browser integration through a client playground;
\item executable functional validation for F1 and F2.
\end{itemize}

\subsection{Out of scope}

\begin{itemize}
\item full SVG specification parity;
\item advanced text shaping and full typography stack;
\item complete filter/mask/animation semantics;
\item complete WebGPU-native backend parity in this phase.
\end{itemize}

\section{Contribution}

The thesis contribution is primarily engineering-focused.
It delivers:

\begin{itemize}
\item a modular conversion prototype between SVG semantics and GPU-ready data;
\item explicit typed contracts between parser, geometry, and render layers;
\item reproducible validation flow for parser and style-normalization functionality;
\item a documented baseline suitable for future optimization and feature expansion.
\end{itemize}

The value is practical: developers gain a clearer path from SVG authoring to explicit rendering workflows.
The value is also methodological: design decisions are linked to geometric and architectural constraints rather than ad hoc implementation.

\section{Thesis structure}

The remainder of the thesis is organized as follows:

\begin{itemize}
\item \textbf{Abstract} summarizes the problem, contribution, and validated functionality.
\item \textbf{Chapter \ref{chap:svg_foundations}} presents SVG and vector-graphics foundations relevant to conversion.
\item \textbf{Chapter \ref{chap:geometry_conversion}} discusses tessellation, triangulation, and robustness strategy.
\item \textbf{Chapter \ref{chap:system_design}} details architecture and implementation of \texttt{Svg2GPU}.
\item \textbf{Chapter \ref{chap:validation}} presents evaluation methodology and executable validation results.
\item \textbf{Chapter \ref{conclusions}} summarizes conclusions and future directions.
\end{itemize}

\section{Closing note}

This thesis treats SVG-to-GPU conversion as a systems problem: semantic input, geometric transformation, explicit rendering, and reproducible verification.
The following chapters develop this perspective from theory to implementation and evaluation.
